%%%%%%%%%%  scan idx 63  |  folio 50  %%%%%%%%%%
\begin{equation}\tag{2.18}
 \gamma(Y)\gamma(Z)=\gamma(Y.Z)\qquad\mathrm{mod.}\ \underline K(X)^{i+j+1}.
\end{equation}

On
fera attention à ce qu'on n'a pas en général égalité dans (2.18). On verra au \S 4 que, si $X$ est \emph{quasi-projective non singulière, la structure d'anneau de $\underline K(X)$ est compatible avec sa filtration et que l'anneau gradué associé $G(\underline K(X))$ s'identifie à un quotient de $A(X)$ par un sous-groupe de torsion (peut-être toujours réduit à $0$\footnote{Non~; pour un contre-exemple cf. ce Séminaire XIV 4.7.})}.

Je ne sais définir directement les $\lambda^i(F)$ en termes de faisceaux que si le corps $k$ est de caractéristique $0$ (c'est la raison qui nous a empêché de démontrer le théorème de Riemann--Roch lorsque $k$ n'est pas de caractéristique $\neq0$). Soit $F$ un faisceau algébrique cohérent sur $X$ sur lequel opère le groupe $\mathfrak S_p$ des permutations de $p$ lettres~; on note alors $(F)^{\mathrm{alt}}$ l'ensemble des $x\in F$ tels que $s.x=\varepsilon(s)x$ pour tout $s\in\mathfrak S_p$ ($\varepsilon(s)$ est la \emph{signature} de $s$). C'est évidemment un sous-faisceau algébrique cohérent de $F$. Ceci dit, on a~:

\textbf{Théorème 2.8.} \emph{Soit $X$ une variété quasi-projective non singulière sur le corps} $k$ de caractéristique $0$. \emph{Soit $F$ un faisceau algébrique cohérent sur $X$~; on a alors~:}
\begin{equation}\tag{2.19}
 \lambda^p(\gamma(F))=\sum_{i\geq0}(-1)^i\gamma\left((\underline{\Tor}_i(F,F,\ldots,F))^{\mathrm{alt}}\right)
\end{equation}

N.B. Il est clair que le groupe $\mathfrak S_p$ agit sur le faisceau $\underline{\Tor}_i(F,\ldots,F)$.

\emph{Principe de la démonstration}~: on considère une résolution finie $\underline L$ de $F$ par des faisceaux localement libres $L_i$, et l'on utilise le lemme général suivant~:

